\section{Background}
\label{sec:background}

\subsection{Standard ReCom}

Let $G = (V, E)$ be the census-tract adjacency graph for a state, with
$|V| = n$ tracts.
A redistricting plan $\pi : V \to [k]$ is valid if every district
$D_i = \pi^{-1}(i)$ is connected in $G$ and population-balanced:
$|P(D_i) - P_{\mathrm{ideal}}| \leq \epsilon \cdot P_{\mathrm{ideal}}$,
where $P_{\mathrm{ideal}} = \sum_{v} \mathrm{pop}(v) / k$ and
$\epsilon = 0.005$.

Standard \recom{}~\citep{deford2021} defines a Markov chain on valid plans
$\pi \in \Pi(G, k)$ as follows.
At each step:
\begin{enumerate}
  \item Select an adjacent district pair $(d_i, d_j)$ uniformly at random
        from all pairs sharing a graph edge.
  \item Let $R = D_i \cup D_j$ (the merged region).
  \item Sample a random spanning tree $T \sim \mathrm{UST}(G[R])$.
  \item Let $C(T)$ be the set of edges $e \in T$ such that removing $e$
        produces a balanced $(A, B)$ bipartition of $R$.
  \item If $C(T) = \emptyset$: reject (stay at $\pi$).
  \item Otherwise: sample $e^* \sim \mathrm{Uniform}(C(T))$; let $\pi'$
        be the plan with $D_i' = A$, $D_j' = B$.
  \item Accept: set $\pi \gets \pi'$.
\end{enumerate}
Steps 5--7 are unconditional: every non-empty balanced cut is accepted.
The chain has no rejection step beyond the empty-cut failure in step 5.

\subsection{The Balanced-Cut Selection Bias}

The transition probability from $\pi$ to $\pi'$ under standard \recom{}
is proportional to the expected number of spanning trees of $G[R]$ that
contain a balanced cut producing $(A, B)$:
\[
  Q_{\mathrm{std}}(\pi \to \pi') \;\propto\;
  \mathbf{E}_{T \sim \mathrm{UST}(G[R])}\bigl[\mathbf{1}[T \text{ cuts to } (A,B)]\bigr].
\]
For the chain to be reversible with respect to the uniform distribution
$\mu_{\mathrm{unif}}$ over $\Pi(G, k)$, we need detailed balance:
$\mu_{\mathrm{unif}}(\pi) \cdot Q(\pi \to \pi') = \mu_{\mathrm{unif}}(\pi') \cdot Q(\pi' \to \pi)$.
Since $\mu_{\mathrm{unif}}(\pi) = \mu_{\mathrm{unif}}(\pi')$, this requires
$Q(\pi \to \pi') = Q(\pi' \to \pi)$---that is, the expected number of trees
containing a $(A,B)$ cut from $\pi$'s side must equal the expected number
containing the same cut from $\pi'$'s side.

In general, this is false.
The merged region $R = D_i \cup D_j$ is the same graph regardless of
direction, so one might expect the forward and reverse transition probabilities
to be equal.
However, the pair-selection probabilities differ: the probability of selecting
the pair $(d_i, d_j)$ may differ between $\pi$ and $\pi'$ because the set of
adjacent district pairs depends on the plan.
Moreover, even fixing the pair, the number of valid balanced cuts per spanning
tree is the same from both sides only if the spanning tree distribution is
symmetric with respect to the $(A, B)$ labelling, which need not hold for
general balanced cuts.

\begin{remark}[Stationarity of standard \recom{}]
\label{rem:std-stationarity}
The stationary distribution of standard \recom{} is not the uniform
distribution over valid plans $\Pi(G,k)$.
It weights plans by their number of spanning tree ``entry paths''---plans
reachable by many distinct spanning tree cuts are overrepresented.
DeFord et al.~\citep{deford2021} do not claim that standard \recom{}
targets the uniform distribution; they use it as a heuristic sampler
whose output is interpreted relative to the chain's own distribution.
For ensemble analysis requiring a specific target distribution, a corrected
chain is needed.
\end{remark}

\subsection{Wilson's Uniform Spanning Tree Algorithm}

Wilson's algorithm~\citep{wilson1996} generates a sample from the uniform
spanning tree distribution $\mathrm{UST}(H)$ of any connected graph $H$ via
loop-erased random walks.
Starting from an arbitrary root $r$, the algorithm constructs the spanning
tree one path at a time: perform a random walk from an unvisited vertex $v$
until hitting a previously visited vertex; erase all loops in the path; add
the resulting loop-erased path to the tree; repeat until all vertices are
covered.

The running time of Wilson's algorithm is $O(|V(H)| \cdot t_{\mathrm{hit}})$,
where $t_{\mathrm{hit}}$ is the maximum hitting time of the random walk on
$H$.
For planar graphs (such as census-tract adjacency graphs), Aldous~\citep{aldous1991}
establishes that $t_{\mathrm{hit}} = O(|V(H)|)$, so the expected running time
is $O(|V(H)|^2)$ in the worst case.
In practice, for the merged district regions $R = D_i \cup D_j$ with
$|R| \approx 2n/k$, Wilson's algorithm runs in $O((n/k)^2)$ expected time---which
is $O(n/k)$ expected steps per vertex, substantially better than the worst-case
bound for non-planar graphs.

For each \fr{} step, Wilson's algorithm is called \emph{twice}: once for
$T_{\mathrm{fwd}}$ (from the current plan's district pair) and once for
$T_{\mathrm{rev}}$ (from the proposed plan's district pair, which is the
same region $R$ but with the balanced-cut constraint evaluated from the other
side).
The doubled call is the direct computational cost of the Metropolis-Hastings
correction.
